\notechapter{N-20221209}{Set-Theoretic and Combinatorial Exercises after the First Lecture}

\section{What this page is doing}

This note is a handwritten exercise continuation of the set-and-counting lecture.  The problems involve floor functions, set-builder notation, binomial identities, lattice or coordinate sets, and counting arguments.  The common goal is to learn how to translate a compact set expression into something countable.

\section{A floor-function example}

One example asks for a set defined by inequalities and floor functions.  The handwritten solution separates the problem into intervals because
\[
  \lfloor x\rfloor=k\quad\Longleftrightarrow\quad k\le x<k+1.
\]
This is the essential method.  Any problem involving \(\lfloor x\rfloor\) should first be split into integer strips.

For instance, if a condition contains
\[
  3\lfloor x\rfloor-5<0,
\]
then
\[
  \lfloor x\rfloor<\frac53,
\]
so
\[
  \lfloor x\rfloor\le1.
\]
One then intersects this with the remaining condition, such as a quadratic inequality.  The final answer is an interval or a finite union of intervals.

\section{A binomial counting identity}

Another exercise invokes the binomial theorem.  The handwritten note uses the idea that terms in a sum can be counted by choosing which positions have a particular property.  A typical identity is
\[
  \sum_{j=0}^n\binom nj=2^n,
\]
or, with signs,
\[
  \sum_{j=0}^n(-1)^j\binom nj=0.
\]

The important thought is not the formula alone.  The coefficient \(\binom nj\) counts ways of choosing \(j\) positions among \(n\), and the binomial theorem packages all choices simultaneously.

\section{Solving a system by discriminants}

One problem gives a two-variable system whose equations can be transformed into a quadratic with a parameter \(p\).  The handwritten solution examines the discriminant
\[
  \Delta\ge0
\]
to decide when real solutions exist.

This is a recurring analytic-algebraic pattern: if a system can be reduced to a quadratic equation, existence of real solutions is controlled by the discriminant.  The solution set changes when the discriminant becomes zero; this is the boundary between two real roots and no real roots.

\section{A prime-number exercise}

Another exercise involves primes \(p,q\).  The handwritten reasoning uses parity and basic divisibility.  The standard approach is:
\begin{enumerate}[(i)]
  \item reduce the equation modulo a small integer such as \(2,3,4\);
  \item use the fact that the only even prime is \(2\);
  \item separate small exceptional cases before treating odd primes.

\end{enumerate}

This is the same habit developed in earlier divisibility notes: do not solve a Diophantine equation blindly.  First use congruences to force the shape of the variables.

\section{A nested-set exercise}

The note includes a Russian sentence asking to ``describe the next set'' or ``write the following set in ordinary language.''  The examples involve finite sets such as
\[
  \varnothing,\qquad \{a_1\},\qquad \{a_1,a_2\},\qquad\ldots
\]
and collections of subsets.  The intended lesson is the difference between an element and a subset:
\[
  a\in A\quad\text{versus}\quad \{a\}\subseteq A.
\]
Confusing these two is one of the most common early set-theory mistakes.

\section{A recurrence of sets}

One exercise defines sets recursively, for example
\[
  A_{k+1}=A_k\cup\{x_{k+1}\}
\]
or by adding all new subsets involving a new element.  The handwritten solution proves the pattern by induction: assume the formula for \(A_k\), then show the construction gives \(A_{k+1}\).

This is the set-theoretic analogue of a recurrence relation.  Instead of numbers changing with \(k\), collections of objects change with \(k\).

\section{An intersection-counting problem}

The final exercise on the second page asks whether a selected element belongs to an intersection of many sets.  The handwritten solution builds sets \(B(A_i)\) and counts their sizes, arriving at a large product-like number.  The exact numerical expression is less important than the logic:
\begin{enumerate}[(i)]
  \item translate membership into a condition;
  \item count how many sets contain a fixed element;
  \item use complements or intersections to isolate the desired elements.

\end{enumerate}
This is a typical finite-set counting method: count incidence pairs in two ways.


\section{Set-builder notation as a translation problem}

A complicated set-builder expression should be read in stages.  First identify the universe of discourse.  Then translate each condition into an inequality, divisibility condition, or membership condition.  Finally decide whether the problem is asking for the set itself, its size, or a proof of equality.

For example, a condition involving \(\lfloor x\rfloor\) is not smooth; it divides the real line into half-open intervals.  A condition involving binomial coefficients often asks for a counting interpretation rather than algebraic simplification.

\section{Recursive set construction}

When a set is defined recursively, the first few stages should be computed explicitly.  This reveals whether the construction is increasing, periodic, stabilizing, or producing genuinely new elements.  After that, induction is the natural proof method.

The general pattern is:
\[
  S_0\subseteq S_1\subseteq S_2\subseteq\cdots
\]
if the rule is monotone.  If the universe is finite, an increasing chain must eventually stabilize.  This is a useful idea behind many finite combinatorial processes.

\section{Expanded synthesis and advanced context}

The exercises are elementary, but they are training the language of discrete mathematics.  Floor functions partition the real line into strips; binomial coefficients count choices; discriminants detect existence; recursive set definitions invite induction; incidence counting turns membership questions into arithmetic.  The high-level skill is always the same: replace notation by a structure that can actually be counted.



\paragraph{Expanded synthesis.}
For this chapter, the elementary material should be read as a study of what remains invariant when coordinates change. The earlier sections (`What this page is doing`, `A floor-function example`, `A binomial counting identity`, `Solving a system by discriminants`, `A prime-number exercise`) compute with arrays, bases, pivots, determinants, or eigenvectors; the deeper object is a linear map together with the spaces it acts on. A matrix is a representation of that map after bases have been chosen. This is why formulas such as
\[
  A' = P^{-1}AP,\qquad \widetilde A = T^{-1}AS
\]
are not cosmetic: they distinguish an intrinsic morphism from a coordinate description.

The structural hierarchy is as follows. Kernels record what a map forgets, images record what it reaches, cokernels record obstructions to solvability, and quotients record information after an equivalence has been imposed. Determinants act on the top exterior power and therefore measure oriented volume; traces are invariant under cyclic rearrangement and become characters in representation theory; adjoints depend on an inner product and lead to self-adjoint spectral theory.

A useful way to return to the computations is to ask, for every formula, which invariant it is measuring. Row reduction measures rank and kernel; diagonalization measures decomposition into invariant subspaces; Gram--Schmidt measures orthogonal projection; positive definiteness measures ellipsoid geometry; tensor notation measures how components transform. The elementary methods are therefore the finite-dimensional laboratory in which duality, quotienting, functoriality, and spectral decomposition can be seen clearly.


\section{Elementary-to-advanced bridge: finite differences and polynomial degree}

Many binomial and floor-sum exercises can be clarified by finite differences.  For a sequence $a(n)$, define
\[
  \Delta a(n)=a(n+1)-a(n).
\]
If $a(n)$ is a polynomial of degree $d$, then $\Delta a(n)$ has degree $d-1$, and
\[
  \Delta^{d+1}a(n)=0.
\]
This is the discrete analogue of differentiating a polynomial until it vanishes.

The binomial coefficient basis is especially natural:
\[
  \Delta \binom nk=\binom n{k-1}.
\]
Therefore integer-valued polynomials are naturally expanded in the basis
\[
  \binom n0,\ \binom n1,\ \binom n2,\ \ldots.
\]

This explains why binomial sums repeatedly appear in discrete counting problems: they are the coordinate system in which finite difference is simplest.

\section{Elementary-to-advanced bridge: polynomial method}

The polynomial method proves combinatorial statements by constructing a polynomial with too many roots or with controlled degree.  The elementary fact is: a nonzero polynomial of degree $d$ over a field has at most $d$ roots.  In several variables, one uses more refined versions involving degree and vanishing on grids.

For example, if a polynomial $P(x)$ of degree at most $n-1$ vanishes at $n$ distinct points, then $P=0$.  Lagrange interpolation gives the exact reconstruction formula:
\[
  P(x)=\sum_{i=1}^n P(a_i)\prod_{j\ne i}\frac{x-a_j}{a_i-a_j}.
\]
Counting information can therefore be encoded into polynomial values.

This is a research-level extension of elementary identities: prove a combinatorial impossibility by forcing a polynomial to be zero while one coefficient or value remains nonzero.

\section{Elementary-to-advanced bridge: Burnside's lemma}

When the note counts objects with symmetry, the advanced principle is group action.  If a finite group $G$ acts on a finite set $X$, the number of orbits is
\[
  |X/G|=\frac1{|G|}\sum_{g\in G}|\operatorname{Fix}(g)|.
\]
This is Burnside's lemma.

It explains why counting up to rotation or reflection requires averaging fixed points rather than simply dividing by the number of symmetries.  Some objects have nontrivial stabilizers, so not every orbit has the same size.

\section{Returning to the original exercises}

Floor sums, binomial identities, and symmetry counts are not separate genres.  Finite differences handle discrete polynomial behavior; the polynomial method turns combinatorics into algebra; group actions handle symmetry.  These are three advanced continuations of the elementary counting toolkit.

% Second-pass deepening for waves 2-4: wave 3 A-C part 1
\section{Second-pass deepening: floor functions as lattice-point counts}

A floor expression such as
\[
  \left\lfloor \frac{an+b}{m}\right\rfloor
\]
counts how many integer thresholds have been crossed.  Sums of floor functions often count lattice points below a line.  For example,
\[
  \sum_{k=0}^{m-1}\left\lfloor \frac{ak}{m}\right\rfloor
\]
counts integer points under a rational-slope diagonal in a rectangle.

This links elementary floor identities to Ehrhart theory, where one counts lattice points in dilates of polytopes:
\[
  L_P(n)=|nP\cap\mathbb Z^d|.
\]
For integral polytopes, $L_P(n)$ is a polynomial in $n$.  Thus floor-sum exercises are one-dimensional shadows of lattice-point enumeration.

\section{Second-pass deepening: group actions as corrected division}

When counting objects up to symmetry, naive division by $|G|$ fails if some objects have symmetry.  Burnside's lemma repairs this by averaging fixed points.  This idea should be tied back to elementary examples: necklaces, colorings of polygons, and arrangements under rotation/reflection.

The orbit-stabilizer theorem explains the correction:
\[
  |Gx|=\frac{|G|}{|\operatorname{Stab}(x)|}.
\]
Objects with larger stabilizers have smaller orbits and must be weighted correctly.  This is why symmetry counting is more subtle than ordinary counting.
